\worksheettitle
  {Измеряем и строим отрезки}
  {\modulemark{M02} Версия учителя}

\begin{teacherbox}[Паспорт модуля]
  \textbf{Результат:} ученик измеряет отрезок, в том числе не от нулевой
  отметки, строит отрезок заданной длины и распознаёт равные отрезки.

  \textbf{Предварительно:} различает отрезок, прямую и луч; знает, что
  $1\text{ см}=10\text{ мм}$. Печатать измерительные страницы в масштабе
  100\%. Точность построения проверяется обычной линейкой ученика.
\end{teacherbox}

\section{Вход и смысл длины}

Во входе: отрезок; линейка; $10$ мм. В заполнении: отрезок —
\textbf{фигура}, длина — \textbf{число} с единицей; запись сообщает
\textbf{длину} отрезка.

\segmentlengthfigure

\begin{teacherbox}[Речевая точность]
  Допустимая бытовая фраза «отрезок равен пяти сантиметрам» уточняется:
  «длина отрезка равна пяти сантиметрам». Это отделяет объект от величины.
\end{teacherbox}

\section{Измерение}

\correctmeasurementfigure

В первом примере начальное показание равно $0$, длина
$5\text{ см }7\text{ мм}$.

\shiftedmeasurementfigure

Во втором примере:
\[
  6\text{ см }8\text{ мм}-1\text{ см }2\text{ мм}
    =5\text{ см }6\text{ мм}.
\]

\readrulerpracticefigure

\begin{practicebox}[Ответ к заданию 1]
  $AB=4\text{ см }6\text{ мм}$. Для второго отрезка:
  \[
    CD=7\text{ см }1\text{ мм}-1\text{ см }3\text{ мм}
      =5\text{ см }8\text{ мм}.
  \]
\end{practicebox}

\begin{teacherbox}[Диагностика]
  Если ученик называет $7\text{ см }1\text{ мм}$, спросите: «Какое
  расстояние прошёл левый конец до начала отрезка?» Сначала возвращайте к
  смыслу разности, а не к заученной команде «вычесть».
\end{teacherbox}

\section{Построение и равенство}

\constructionalgorithmfigure
\studentconstructionanswersfigure

\begin{teacherbox}[Проверка задания 2]
  $AB=5\text{ см }8\text{ мм}$, $CD=3\text{ см }5\text{ мм}$. Проверяйте
  положение обоих концов и длину, а не совпадение с напечатанным образцом.
\end{teacherbox}

\equalcontrastingsegmentsfigure

\begin{practicebox}[Ответ к заданию 3]
  Равны $AB$ и $CD$: обе длины равны $4\text{ см}$. Отрезок $MN$
  имеет длину $6\text{ см }6\text{ мм}$.
\end{practicebox}

\section{Ошибка и контроль}

\begin{warningbox}[Исправление]
  Причина: показание правого конца принято за длину; начальный участок линейки
  не вычтен. Верно:
  \[
    7\text{ см }1\text{ мм}-1\text{ см }3\text{ мм}
      =5\text{ см }8\text{ мм}.
  \]
\end{warningbox}

\begin{checkpointbox}[Ответы]
  \begin{enumerate}[label=\textbf{\arabic*.}]
    \item $8\text{ см }1\text{ мм}-2\text{ см }4\text{ мм}
      =5\text{ см }7\text{ мм}$.
    \item Построение $MN=4\text{ см }6\text{ мм}$ с двумя отмеченными
      концами.
    \item Да. Равные длины означают равные отрезки независимо от положения.
  \end{enumerate}
\end{checkpointbox}

\begin{teacherbox}[Решение о маршруте]
  M02-R назначается при подмене длины показанием конца или отсутствии разности
  при сдвинутом начале. M02-H — при верном способе и вычислительной или
  построительной неточности. Устойчивый контроль открывает M03.
\end{teacherbox}
